\chapter{Python: ambiente e fundamentos}
\label{cap:python-fundamentos}

\section{O papel do Python neste livro}

Python é uma linguagem de programação de propósito geral. Neste livro, ela será usada como ferramenta de investigação: realizaremos cálculos, analisaremos dados, construiremos gráficos, desenharemos vetores e simularemos modelos físicos.

O código não substitui a formulação matemática nem a interpretação física. Ele permite executar um procedimento de forma reproduzível e explorar situações que seriam trabalhosas manualmente. Um resultado numérico só se torna fisicamente útil depois que suas unidades, hipóteses e limites são examinados.

Adotaremos o \textbf{Python 3.12}. Versões posteriores da série 3 provavelmente executarão a maior parte dos exemplos, mas a versão 3.12 será nossa referência de desenvolvimento e validação. Fixar uma versão reduz diferenças de comportamento entre computadores e facilita a reprodução dos resultados.

\section{Interpretador, programa e script}

O interpretador Python é o programa que lê e executa instruções escritas na linguagem Python. Uma instrução pode ser digitada em uma sessão interativa ou armazenada em um arquivo de texto com extensão \texttt{.py}.

Um arquivo como

\begin{center}
  \texttt{primeiro\_programa.py}
\end{center}

é chamado de \emph{script}. No terminal integrado do VS Code, ele pode ser executado com

\begin{lstlisting}[language={},numbers=none]
python primeiro_programa.py
\end{lstlisting}

Neste projeto, o interpretador do ambiente virtual pode ser selecionado diretamente pelo VS Code. Isso garante que o editor utilize a mesma versão de Python e as mesmas bibliotecas usadas na validação do livro.

\section{Onde escrever e executar Python}

Python não está vinculado a um único programa. Um arquivo com extensão \texttt{.py} pode ser escrito em diferentes editores e ambientes, desde que um interpretador compatível e as bibliotecas necessárias estejam disponíveis.

É útil distinguir três conceitos:

\begin{itemize}
  \item um \textbf{editor de código} oferece recursos para escrever e organizar arquivos;
  \item uma \textbf{IDE}, ou ambiente integrado de desenvolvimento, reúne editor, execução, depuração, terminal e ferramentas de projeto;
  \item um \textbf{notebook} organiza texto, equações, código e resultados em células executáveis.
\end{itemize}

Entre as alternativas mais conhecidas estão:

\begin{description}
  \item[Visual Studio Code:] editor extensível que, com a extensão Python, oferece execução, depuração, terminal integrado e suporte a notebooks. Será o ambiente principal de validação deste livro.

  \item[JupyterLab:] ambiente local orientado a notebooks. É especialmente útil para exploração de dados, testes rápidos e combinação de explicações com resultados computacionais.

  \item[Google Colab:] serviço executado no navegador e baseado em notebooks. Permite escrever e executar Python sem instalar previamente um ambiente completo no computador. A sessão é remota e temporária; arquivos e bibliotecas instalados durante uma sessão podem não permanecer disponíveis na sessão seguinte.

  \item[PyCharm:] IDE dedicada ao desenvolvimento em Python, com ferramentas de navegação, testes e depuração de projetos.

  \item[Spyder:] IDE voltada à computação científica, com editor, console interativo e visualização de variáveis.
\end{description}

Também é possível executar Python diretamente em um terminal ou utilizar outros editores. A escolha do ambiente não altera as leis físicas nem a sintaxe fundamental da linguagem, mas pode mudar a forma de executar arquivos, instalar bibliotecas e visualizar gráficos.

Os exemplos deste livro serão escritos como scripts Python portáveis. Sempre que um recurso depender especificamente de notebook, VS Code ou Colab, essa dependência será indicada. Para garantir que os resultados publicados possam ser reproduzidos, a validação oficial será realizada no ambiente virtual do projeto com Python 3.12.

\section{O primeiro programa}

A função \texttt{print} envia uma representação de um valor para a saída do programa:

\begin{lstlisting}
print("Fisica com Python")
\end{lstlisting}

O texto entre aspas é uma \emph{string}, isto é, uma sequência de caracteres. Ao executar a instrução, o terminal apresenta

\begin{lstlisting}[language={},numbers=none]
Fisica com Python
\end{lstlisting}

Os parênteses delimitam os argumentos. Aqui, o argumento é uma string.

\section{Valores, variáveis e atribuição}

Uma variável é um nome associado a um valor. A instrução

\begin{lstlisting}
massa = 2.5
\end{lstlisting}

associa o nome \texttt{massa} ao valor numérico \texttt{2.5}. O símbolo \texttt{=} realiza uma \textbf{atribuição}; ele não afirma uma igualdade algébrica permanente.

Depois de

\begin{lstlisting}
velocidade = 3.0
velocidade = 5.0
\end{lstlisting}

o nome \texttt{velocidade} está associado a \texttt{5.0}. A segunda atribuição substitui a referência anterior.

Nomes de variáveis devem comunicar seu significado. Usaremos, por exemplo,

\begin{lstlisting}
tempo_total = 12.0
velocidade_inicial = 4.5
\end{lstlisting}

em vez de nomes genéricos como \texttt{a} ou \texttt{x1}, exceto quando a notação matemática justificar nomes curtos.

O Python não associa unidades físicas automaticamente a números comuns. Por isso, as unidades devem ser declaradas no texto, nos comentários e nos rótulos das saídas:

\begin{lstlisting}
distancia = 120.0  # metros
tempo = 10.0       # segundos
\end{lstlisting}

O caractere \texttt{\#} inicia um comentário. O interpretador ignora o restante daquela linha, mas o comentário ajuda uma pessoa a compreender o programa.

\section{Tipos fundamentais}

Cada valor possui um tipo. Entre os tipos fundamentais estão:

\begin{center}
\begin{tabular}{lll}
\toprule
Tipo & Significado & Exemplo\\
\midrule
\texttt{int} & número inteiro & \texttt{12}\\
\texttt{float} & número de ponto flutuante & \texttt{9.81}\\
\texttt{str} & sequência de caracteres & \texttt{tempo}\\
\texttt{bool} & valor lógico & \texttt{True}, \texttt{False}\\
\bottomrule
\end{tabular}
\end{center}

A função \texttt{type} informa o tipo de um valor:

\begin{lstlisting}
print(type(12))
print(type(9.81))
\end{lstlisting}

A saída esperada é

\begin{lstlisting}[style=saidaesperada]
<class 'int'>
<class 'float'>
\end{lstlisting}

Um \texttt{float} é uma aproximação computacional de um número real. Como sua representação possui precisão finita, algumas operações produzem pequenas diferenças de arredondamento. Esse comportamento será considerado sempre que compararmos resultados numéricos.

\section{Operadores aritméticos}

Os operadores aritméticos básicos são:

\begin{center}
\begin{tabular}{lll}
\toprule
Operador & Operação & Exemplo\\
\midrule
\texttt{+} & adição & \texttt{a + b}\\
\texttt{-} & subtração & \texttt{a - b}\\
\texttt{*} & multiplicação & \texttt{a * b}\\
\texttt{/} & divisão & \texttt{a / b}\\
\texttt{**} & potência & \texttt{a ** 2}\\
\texttt{//} & divisão inteira & \texttt{a // b}\\
\texttt{\%} & resto da divisão & \texttt{a \% b}\\
\bottomrule
\end{tabular}
\end{center}

A precedência das operações é semelhante à utilizada na álgebra. Parênteses devem ser usados quando tornam a intenção mais clara:

\begin{lstlisting}
energia = 0.5 * massa * velocidade**2
media = (medida_1 + medida_2) / 2
\end{lstlisting}

\section{Apresentação de resultados}

Uma \emph{f-string} permite inserir valores em um texto:

\begin{lstlisting}
energia = 9.0
print(f"Energia cinetica: {energia:.2f} J")
\end{lstlisting}

A especificação \texttt{:.2f} apresenta o número com duas casas decimais. A saída será

\begin{lstlisting}[language={},numbers=none]
Energia cinetica: 9.00 J
\end{lstlisting}

Formatar um valor não aumenta a precisão do cálculo. A formatação controla apenas como o número é exibido.

\section{Funções}

Uma função reúne instruções associadas a uma tarefa. Ela pode receber valores de entrada, chamados argumentos, e devolver um resultado:

\begin{lstlisting}
def energia_cinetica(massa, velocidade):
    return 0.5 * massa * velocidade**2
\end{lstlisting}

A palavra \texttt{def} inicia a definição; os dois-pontos abrem um bloco; a indentação identifica as instruções pertencentes à função; e \texttt{return} devolve o resultado.

A função pode ser utilizada por meio de uma chamada:

\begin{lstlisting}
energia = energia_cinetica(2.0, 3.0)
print(energia)
\end{lstlisting}

A saída esperada é

\begin{lstlisting}[style=saidaesperada]
9.0
\end{lstlisting}

Separar um cálculo em uma função reduz repetição, facilita testes e explicita o modelo empregado.

\section{Módulos e bibliotecas}

Um módulo é um arquivo que contém definições Python reutilizáveis. Uma biblioteca é uma coleção organizada de módulos. O comando \texttt{import} torna essas definições disponíveis no programa.

O módulo \texttt{math} pertence à biblioteca padrão, instalada juntamente com o Python:

\begin{lstlisting}
import math

raiz = math.sqrt(25)
angulo = math.radians(30)
\end{lstlisting}

O prefixo \texttt{math.} informa em qual módulo a função foi definida e evita ambiguidades.

Bibliotecas externas são instaladas separadamente. Neste livro, as principais serão:

\begin{itemize}
  \item \textbf{NumPy}: arrays e computação numérica;
  \item \textbf{Matplotlib}: gráficos e visualizações;
  \item \textbf{SciPy}: algoritmos científicos;
  \item \textbf{SymPy}: matemática simbólica.
\end{itemize}

Um exemplo de importação é

\begin{lstlisting}
import numpy as np
import matplotlib.pyplot as plt
\end{lstlisting}

Os nomes \texttt{np} e \texttt{plt} são convenções amplamente utilizadas. Eles são aliases: nomes abreviados associados aos módulos importados.

\section{Arquivos, pastas e caminhos com pathlib}
\label{sec:pathlib}

Os programas que produzem gráficos precisam informar onde os arquivos serão guardados. Um caminho completo preso a uma pasta pessoal faria o código funcionar apenas naquela organização. Para criar caminhos portáteis, usaremos a classe \texttt{Path}, pertencente ao módulo \texttt{pathlib} da biblioteca padrão do Python.

\begin{lstlisting}
from pathlib import Path

pasta_atual = Path.cwd()
print(pasta_atual)
\end{lstlisting}

O método \texttt{Path.cwd()} informa a \textbf{pasta de trabalho atual}: o local a partir do qual o programa está sendo executado. A saída depende do computador e da pasta aberta no terminal. No Windows, ela poderia ser

\begin{lstlisting}[style=saidaesperada]
C:\Users\ana\FisicaComPython\livro
\end{lstlisting}

No Linux ou no Google Colab, o mesmo comando pode produzir algo como

\begin{lstlisting}[style=saidaesperada]
/home/ana/FisicaComPython/livro
\end{lstlisting}

Essas diferenças não exigem alterações no restante do programa. Em vez de escrever manualmente barras próprias de cada sistema, o operador \texttt{/} combina partes de um caminho:

\begin{lstlisting}
pasta_figuras = Path("figuras") / "cap00"
arquivo = pasta_figuras / "meu-grafico.pdf"
\end{lstlisting}

Como o caminho começa com \texttt{figuras} e não com uma unidade de disco ou uma barra inicial, ele é \textbf{relativo}. Se a pasta de trabalho for \texttt{livro}, o arquivo ficará em \texttt{livro/figuras/cap00/meu-grafico.pdf}. O método \texttt{resolve} permite visualizar o caminho completo calculado no computador:

\begin{lstlisting}
print(arquivo.resolve())
\end{lstlisting}

Antes de salvar, precisamos garantir que a pasta exista:

\begin{lstlisting}
pasta_figuras.mkdir(parents=True, exist_ok=True)
\end{lstlisting}

O argumento \texttt{parents=True} autoriza a criação de pastas intermediárias; \texttt{exist\_ok=True} evita um erro quando a pasta já existe. A sequência portátil que adotaremos será, portanto,

\begin{lstlisting}
from pathlib import Path

pasta_figuras = Path("figuras") / "cap00"
pasta_figuras.mkdir(parents=True, exist_ok=True)
arquivo = pasta_figuras / "meu-grafico.pdf"
\end{lstlisting}

Nos códigos oficiais deste livro, o terminal será aberto na pasta \texttt{livro}. Por exemplo,

\begin{lstlisting}[language={},numbers=none]
python codigos/cap00/primeiro_grafico.py
\end{lstlisting}

Assim, os programas guardam os resultados na pasta \texttt{figuras}, dentro de \texttt{livro}.

Se o leitor copiar um script para outro projeto, a mesma instrução criará a pasta \texttt{figuras} no novo local de execução.

Também existe a variável especial \lstinline|__file__|, que informa o arquivo do script. Ela pode ser útil em projetos com estruturas rígidas, mas não está disponível da mesma forma em todos os notebooks. Por isso, os exemplos introdutórios usarão caminhos relativos à pasta de trabalho. No VS Code, no JupyterLab e no Colab, o leitor deve conferir \texttt{Path.cwd()} quando tiver dúvida sobre o destino de um arquivo.

\section{Primeiros gráficos com Matplotlib}
\label{sec:primeiros-graficos}

Um gráfico científico não é apenas uma ilustração: ele representa a relação entre grandezas. Antes de interpretar curvas de posição, velocidade, energia ou campo, precisamos compreender como os dados são convertidos em elementos visuais e como esses elementos são construídos em Python.

Usaremos o módulo \texttt{matplotlib.pyplot}, importado com o alias \texttt{plt}. A biblioteca Matplotlib organiza um gráfico em dois objetos principais:

\begin{description}
  \item[Figura (\texttt{fig}):] a área completa que será exibida ou salva, como uma folha que pode conter um ou mais gráficos;
  \item[Eixos (\texttt{ax}):] a região em que os dados são desenhados, incluindo os eixos coordenados, marcas, rótulos e curvas.
\end{description}

Essa distinção explica a instrução

\begin{lstlisting}
fig, ax = plt.subplots()
\end{lstlisting}

A função \texttt{plt.subplots} cria a figura e um sistema de eixos. Os nomes \texttt{fig} e \texttt{ax} são convenções, não palavras reservadas do Python.

\subsection{Um primeiro gráfico a partir de dados}

Considere posições medidas a cada segundo. Duas listas armazenam os valores correspondentes de tempo e posição. O primeiro elemento de \texttt{tempo} é associado ao primeiro elemento de \texttt{posicao}, e assim sucessivamente.

\lstinputlisting[
  caption={Primeiro gráfico construído a partir de dados de posição e tempo.},
  label={cod:primeiro-grafico}
]{codigos/cap00/primeiro_grafico.py}

A saída textual esperada é

\lstinputlisting[
  style=saidaesperada,
  caption={Saída esperada do primeiro programa gráfico.}
]{codigos/cap00/saidas/primeiro_grafico.txt}

O arquivo produzido pelo programa é mostrado na \cref{fig:primeiro-grafico}.

\begin{figure}[htbp]
  \centering
  \includegraphics[width=0.78\textwidth]{cap00/primeiro-grafico.pdf}
  \caption{Posição em função do tempo. Os círculos identificam os dados fornecidos ao programa e os segmentos apenas orientam a leitura entre eles.}
  \label{fig:primeiro-grafico}
\end{figure}

No \cref{cod:primeiro-grafico}:

\begin{itemize}
  \item \texttt{ax.plot(tempo, posicao)} associa os valores da primeira sequência ao eixo horizontal e os da segunda ao eixo vertical;
  \item o argumento \texttt{marker="o"} marca cada par de dados com um círculo;
  \item \texttt{set\_xlabel} e \texttt{set\_ylabel} identificam as grandezas e suas unidades;
  \item \texttt{set\_title} acrescenta um título descritivo;
  \item \texttt{tight\_layout} ajusta os espaços para evitar cortes nos rótulos;
  \item \texttt{savefig} salva a figura, enquanto \texttt{plt.close} libera a memória usada por ela.
\end{itemize}

Durante uma exploração interativa, \texttt{plt.show()} pode ser usado para abrir a figura em uma janela ou exibi-la em um notebook. Nos scripts do livro, preferimos \texttt{savefig}: o resultado fica registrado, pode ser reproduzido e pode ser inserido no texto. O formato PDF preserva linhas e letras como elementos vetoriais, adequados à impressão.

\subsection{Amostragem de um intervalo com NumPy}

Listas são suficientes para armazenar uma pequena quantidade de dados experimentais. Para representar uma curva, calculamos a função em muitos pontos.

Por exemplo,

\begin{lstlisting}
tempo = np.linspace(0, 5, 101)
\end{lstlisting}

cria um array NumPy com 101 valores igualmente espaçados, incluindo \num{0} e \num{5}. Um \emph{array} é uma estrutura numérica que permite aplicar uma operação a todos os seus elementos. Assim,

\begin{lstlisting}
posicao = 2.0 * tempo
\end{lstlisting}

calcula uma posição para cada instante sem escrever um laço explicitamente. O conjunto de valores escolhidos para a variável independente é o \textbf{domínio amostrado}. A curva exibida é formada por segmentos entre os pontos calculados; ela não contém infinitos pontos. Aumentar a amostragem pode tornar a curva mais suave, mas não corrige um modelo físico inadequado.

\subsection{Eixos, estilos, legenda, grade e anotação}

Quando duas ou mais curvas aparecem nos mesmos eixos, cada uma precisa ser identificável também em impressão sem cores. Para isso, combinamos rótulos com estilos de linha diferentes.

\lstinputlisting[
  caption={Elementos essenciais de um gráfico científico.},
  label={cod:elementos-grafico}
]{codigos/cap00/elementos_grafico.py}

A saída esperada é

\lstinputlisting[
  style=saidaesperada,
  caption={Saída esperada do programa com duas curvas.}
]{codigos/cap00/saidas/elementos_grafico.txt}

\begin{figure}[htbp]
  \centering
  \includegraphics[width=0.82\textwidth]{cap00/elementos-grafico.pdf}
  \caption{Dois movimentos apresentados nos mesmos eixos. Cor, estilo de linha, legenda e anotação trabalham juntos para comunicar a informação.}
  \label{fig:elementos-grafico}
\end{figure}

Os novos comandos controlam diferentes partes da comunicação visual:

\begin{description}
  \item[\texttt{color}:] escolhe a cor da curva;
  \item[\texttt{linestyle}:] define o estilo, como linha contínua ou tracejada;
  \item[\texttt{label}:] associa um texto à curva;
  \item[\texttt{legend}:] exibe os textos associados às curvas;
  \item[\texttt{set\_xlim} e \texttt{set\_ylim}:] definem os limites visíveis dos eixos;
  \item[\texttt{grid}:] acrescenta linhas de referência; o parâmetro \texttt{alpha} controla sua transparência;
  \item[\texttt{annotate}:] relaciona uma observação a um ponto do gráfico. O argumento \texttt{xy} indica o ponto e \texttt{xytext} indica a posição do texto.
\end{description}

Limites de eixos devem ser escolhidos com critério. Cortar parte da escala pode destacar uma variação pequena, mas também pode exagerar visualmente diferenças. Uma grade deve ajudar na leitura, não competir com os dados. Uma legenda é necessária quando há várias curvas sem identificação direta; em uma curva única, bons rótulos de eixos e uma legenda de figura geralmente bastam.

\subsection{Mais de um sistema de eixos}

A função \texttt{plt.subplots} também cria vários gráficos na mesma figura:

\begin{lstlisting}
fig, eixos = plt.subplots(1, 2)
eixos[0].plot(x, y_1)
eixos[1].plot(x, y_2)
\end{lstlisting}

Os argumentos \texttt{1, 2} especificam uma linha e duas colunas. Nesse caso, \texttt{eixos} é uma coleção: \texttt{eixos[0]} seleciona o primeiro sistema de eixos e \texttt{eixos[1]}, o segundo. Esse recurso será usado quando a comparação lado a lado for mais clara do que sobrepor muitas curvas.

\subsection{Linhas de referência, marcas, proporção e formas}

Alguns gráficos exigem mais do que curvas. Os comandos a seguir serão úteis na trigonometria e na representação de vetores:

\begin{lstlisting}
ax.axhline(0)                  # linha horizontal em y = 0
ax.axvline(0)                  # linha vertical em x = 0
ax.set_xticks([0, 1, 2])       # posições das marcas
ax.set_aspect("equal")         # mesma escala visual nos eixos
circulo = plt.Circle((0, 0), 1, fill=False)
ax.add_patch(circulo)          # adiciona a forma aos eixos
ax.quiver(0, 0, 1, 1)         # seta de (0, 0) até (1, 1)
\end{lstlisting}

Os métodos \texttt{axhline} e \texttt{axvline} criam linhas de referência paralelas aos eixos. O método \texttt{set\_xticks} controla onde aparecem as marcas do eixo horizontal; um segundo argumento opcional pode fornecer textos personalizados, como $\pi/2$ e $\pi$. O ajuste \texttt{set\_aspect("equal")} faz uma unidade horizontal ocupar o mesmo comprimento visual que uma unidade vertical, condição necessária para que um círculo não pareça uma elipse.

Uma \texttt{Circle} é uma forma geométrica que só passa a pertencer ao gráfico depois de ser entregue a \texttt{add\_patch}. O método \texttt{quiver} desenha setas e será estudado com mais profundidade quando representarmos vetores. Por enquanto, basta observar que os dois primeiros números indicam a origem e os dois últimos, as componentes da seta.

Textos matemáticos podem ser escritos entre sinais de dólar. Usaremos frequentemente strings como \lstinline|r"$\sin(\theta)$"|. O prefixo \texttt{r} cria uma \emph{raw string}, na qual as barras invertidas são preservadas para a notação matemática do Matplotlib.

\subsection{Critérios para um gráfico científico}

Nem todo cálculo precisa de um gráfico. Ele é útil quando a forma, a tendência, a comparação ou a variação entre grandezas acrescenta informação. Quando for usado, verificaremos se:

\begin{enumerate}
  \item os eixos identificam grandezas e unidades;
  \item o domínio e os limites correspondem ao problema físico;
  \item curvas diferentes podem ser distinguidas sem depender apenas de cor;
  \item a legenda existe somente quando necessária;
  \item a quantidade de pontos representa adequadamente a função ou os dados;
  \item título, grade e anotações ajudam a leitura sem criar excesso visual;
  \item o arquivo foi salvo com resolução ou formato apropriado.
\end{enumerate}

\section{Ambiente virtual e dependências}

Um ambiente virtual isola o interpretador e as bibliotecas de um projeto. Assim, projetos diferentes podem utilizar versões diferentes de uma mesma biblioteca sem interferência entre si.

Neste repositório, o ambiente local fica na pasta \texttt{.venv}. No terminal do VS Code para Windows, ele pode ser ativado com

\begin{lstlisting}[language={},numbers=none]
.\.venv\Scripts\Activate.ps1
\end{lstlisting}

As dependências declaradas no arquivo \texttt{requirements.txt} podem ser instaladas com

\begin{lstlisting}[language={},numbers=none]
python -m pip install -r requirements.txt
\end{lstlisting}

O comando \texttt{python -m pip} garante que o instalador utilizado pertence ao interpretador Python selecionado.

\section{Erros como informação}

Uma mensagem de erro não significa apenas que o programa falhou. Ela registra onde ocorreu o problema e qual regra foi violada. Três categorias aparecerão com frequência:

\begin{itemize}
  \item \texttt{SyntaxError}: a instrução não segue a sintaxe da linguagem;
  \item \texttt{NameError}: um nome foi usado sem estar definido;
  \item \texttt{TypeError}: uma operação recebeu um tipo incompatível.
\end{itemize}

Ao investigar um erro, devemos ler a última linha da mensagem e localizar o arquivo e a linha indicados. Corrigir erros faz parte do processo normal de modelagem computacional.

\section{Critério de validação}

Antes de considerar um código aprovado, verificaremos:

\begin{enumerate}
  \item se o programa executa sem erros;
  \item se as unidades utilizadas são consistentes;
  \item se o resultado coincide com um caso calculado manualmente;
  \item se casos-limite produzem comportamento coerente;
  \item se nomes, comentários e saídas são compreensíveis;
  \item se gráficos possuem eixos, unidades e legendas adequados.
\end{enumerate}

Esses critérios acompanharão todos os códigos do livro.
